% Section 6: Policy Implications
% Word target: ~1000 words

\section{Policy Implications}\label{sec:policy}

The numerical experiments in Section~\ref{sec:experiments} provide
evidence-based guidance for DRT deployment decisions in Chinese cities.
Bidirectional meeting point systems offer substantial vehicle efficiency
gains, but their viability depends critically on walking tolerance, fleet
supply, and demand density. This section translates the experimental
findings into five actionable recommendations for DRT operators and urban
planners, framed for the low-density suburban contexts where DRT is most
commonly deployed in China.

% -----------------------------------------------------------------------
\subsection{R1: Set Walking Radius Threshold at 1000\,m}
% -----------------------------------------------------------------------

\paragraph{Evidence.}
The sensitivity analysis (Section~\ref{sec:experiments}) shows that
FullModel acceptance rate is \textbf{0.0\%} at $\rho \leq 400$\,m and
rises to \textbf{9.0\%} at $\rho = 1000$\,m. Below the 1000\,m threshold,
the MNL utility penalty for walking is so large that no meeting-point offer
clears the outside option, and the bidirectional system delivers no
efficiency benefit over door-to-door service. At $\rho = 1000$\,m,
FullModel vehicle-km drops to 219.7 versus 1693.8 for DoorToDoor in the
standard density tier --- a 87.0\% reduction in vehicle travel.

\paragraph{Policy implication.}
DRT operators in Chinese cities should conduct walking tolerance surveys
before deploying bidirectional meeting points. Districts where residents
routinely walk more than 1000\,m to transit stops --- common in suburban
areas with sparse bus networks --- are prime candidates. Districts with
elderly-dominant populations or poor pedestrian infrastructure (inadequate
lighting, missing footpaths) should retain door-to-door service. Operators
should pilot bidirectional meeting points with opt-in enrollment, allowing
passengers to self-select based on their own walking tolerance before
system-wide rollout. The 1000\,m threshold should be treated as a
\emph{minimum} viable radius; higher thresholds will yield higher acceptance
rates and greater efficiency gains.

% -----------------------------------------------------------------------
\subsection{R2: Minimum Fleet Ratio of 15 Vehicles per 100 Daily Requests}
% -----------------------------------------------------------------------

\paragraph{Evidence.}
Fleet size sensitivity analysis shows that FullModel acceptance peaks at
\textbf{24.0\%} with $n_{\text{vehicles}} = 30$ for 200 requests (15.0
vehicles per 100 requests). Diminishing returns set in beyond 20 vehicles:
the marginal acceptance gain per additional vehicle drops sharply after this
point. At $n_{\text{vehicles}} = 15$ (baseline), DoorToDoor acceptance
exceeds FullModel by \textbf{38 percentage points}, confirming that
door-to-door service is more robust under vehicle scarcity because it does
not impose walking penalties on passengers.

\paragraph{Policy implication.}
Chinese city DRT operators planning new services should budget for at least
15 vehicles per 100 peak-hour requests. Underfleeting negates the efficiency
gains from bidirectional meeting point optimization: a system with too few
vehicles will achieve lower acceptance than a simple door-to-door service,
while still imposing walking burdens on passengers. Operators can phase in
vehicles incrementally, using the sensitivity curve (Figure~\ref{fig:sensitivity})
to project service quality at each fleet size. Bidirectional meeting points
deliver their greatest advantage when fleet supply is sufficient to serve
the majority of requests; the 15-per-100 ratio is the minimum threshold for
this condition to hold.

% -----------------------------------------------------------------------
\subsection{R3: Monitor Time-Sensitive Passengers for Service Equity}
% -----------------------------------------------------------------------

\paragraph{Evidence.}
Equity analysis across three passenger types shows that time-sensitive
passengers have the lowest acceptance rate (\textbf{14.4\%}), compared to
price-sensitive (\textbf{25.4\%}) and walk-sensitive (\textbf{20.2\%})
passengers. Time-sensitive passengers face longer in-vehicle travel times
(avg IVT 1109.6\,s) and higher waiting times (avg 603.2\,s) relative to
other types, making meeting-point offers less attractive under their utility
function. The Gini coefficient of per-type acceptance rates is
\textbf{0.122}, indicating moderate but policy-relevant inequality.

\paragraph{Policy implication.}
DRT operators should track per-type acceptance rates in real deployments.
If time-sensitive passengers are disproportionately rejected, operators
should consider: (a) adjusting MNL utility weights to reduce the penalty
for this type; (b) offering a door-to-door fallback option for passengers
who opt out of meeting-point service; or (c) subsidizing the disadvantaged
type through reduced fares or priority dispatch. Regulatory bodies in
Chinese cities should require DRT operators to report disaggregated
acceptance rates by passenger segment as a condition of operating licenses.
This aligns with TR Part A's equity mandate for public transit policy and
provides a measurable accountability mechanism for service providers.

% -----------------------------------------------------------------------
\subsection{R4: Match Service Mode to City Density Tier}
% -----------------------------------------------------------------------

\paragraph{Evidence.}
The standard scenario (200 requests, 15 vehicles) shows FullModel
vehicle-km = 628.5 versus DoorToDoor = 2603.7, a \textbf{75.9\%} efficiency
gain. However, this gain is density-dependent: at low demand densities
(100 requests per $20\,\text{km}^2$), FullModel acceptance approaches 0\%
at $\rho = 500$\,m, while DoorToDoor remains viable. The three-tier
framework below synthesizes these findings into a practical deployment guide.

\paragraph{Policy implication.}
A three-tier deployment framework for Chinese cities:
\begin{itemize}
  \item \textbf{Tier 1} ($>$300 requests/day in service zone): Bidirectional
        meeting points with full MNL choice model. Demand density is
        sufficient to generate spatially compatible request clusters, and
        the efficiency gains (up to 75.9\% vehicle-km reduction) justify
        the system complexity and passenger walking burden.
  \item \textbf{Tier 2} (100--300 requests/day): Single-sided pickup meeting
        points with door-to-door dropoff. This hybrid mode captures
        approximately half the vehicle-km savings of full bidirectional
        assignment while reducing the walking burden to pickup only.
  \item \textbf{Tier 3} ($<$100 requests/day): Door-to-door service.
        Meeting point overhead exceeds efficiency gains at low demand
        levels; the system cannot generate enough spatially compatible
        clusters to justify passenger walking.
\end{itemize}

\begin{figure}[h]
\centering
[FIGURE 6 PLACEHOLDER: Policy benefit map --- demand density $\times$ walking tolerance]
\caption{Policy deployment map: recommended service mode by demand density and walking tolerance.}
\label{fig:policy-map}
\end{figure}

% -----------------------------------------------------------------------
\subsection{R5: Deploy Rolling Horizon Re-optimization for Dynamic Demand}
% -----------------------------------------------------------------------

\paragraph{Evidence.}
The AblationNoRollingHorizon experiment shows a \textbf{16.6\%} increase in
vehicle-km when rolling horizon re-optimization is disabled (753.4 vs.\
628.5\,vkm). More critically, average waiting time rises from 450.3\,s
(FullModel) to \textbf{3948.8\,s} without rolling horizon --- an 8.8-fold
increase. Static greedy assignment cannot correct early routing decisions
as new requests arrive, leading to cascading delays and inefficient vehicle
positioning.

\paragraph{Policy implication.}
DRT operators should implement rolling horizon re-optimization with a
planning window $H \geq 30$ minutes and a re-optimization time step
$\Delta = 5$ minutes. The computational overhead is acceptable: average
decision time per request remains below 1 second on instances with 200
requests and 15 vehicles, making real-time deployment feasible on standard
server hardware. Chinese city operators should plan for cloud-based dispatch
systems rather than on-vehicle computation, enabling centralized
re-optimization across the full fleet. Static assignment (greedy insertion
without re-optimization) is insufficient for dynamic urban demand patterns
and should not be used in production deployments of bidirectional DRT.

\paragraph{Deployment summary.}
Operators should assess three factors before choosing a service mode:
(1) walking tolerance of the target population ($\rho \geq 1000$\,m
required for bidirectional service); (2) fleet supply ($\geq$15 vehicles
per 100 peak-hour requests); and (3) demand density ($>$300 requests/day
for full bidirectional, 100--300 for single-sided, $<$100 for
door-to-door). The three-tier framework in R4 and the deployment map in
Figure~\ref{fig:policy-map} provide a practical decision tool for Chinese
city DRT planners navigating these trade-offs.
